\section{Fixed particle numbers and the Thomas--Fermi coefficient}
\label{sec:fixed-sectors}

We first recover fixed-sector energy differences from the priced deficit
proved in \cref{sec:limit}.  This step uses monotonicity of the sector
energies, without requiring convexity or requiring every sector to minimize
at some price.  The priced minimum is concave in the price, and its
minimizing sectors bound its secant slopes even when the minimizing sector
changes.  Integrating the deficit law therefore determines the energy at
those sectors; monotonicity then brackets the desired particle number.
We identify the coefficient by the corresponding Thomas--Fermi calculation.

\subsection{Integrating the price and bracketing the sector}

\begin{proposition}\label{prop:fixed-sector-limit}
Let $q$ be the deficit coefficient in \cref{thm:priced-deficit}.  For every
sequence of integers $Z>m\ge1$ such that $m\to\infty$ and $Z/m\to\infty$,
\begin{equation}\label{eq:fixed-sector-limit}
 \frac{E_Z(Z-m)-E_Z(Z)}{m^{7/3}}
       \longrightarrow\frac37q^{-4/3}.
\end{equation}
\end{proposition}

\begin{proof}
Fix such a sequence.  At each fixed $t>0$, choose any index $n_t$ minimizing
$E_n+tm^{4/3}n$, and put
\[
 d_{Z,m}(t)=\frac{Z-n_t}{m},\qquad
 G_{Z,m}(t)=m^{-7/3}
       \bigl[P(tm^{4/3})-tm^{4/3}Z-E_Z(Z)\bigr].
\]
The choice $s=t^{-1/4}m^{-1/3}$ satisfies
\cref{eq:scaling-regime}, since $Zs^3=t^{-3/4}Z/m\to\infty$.
The deficit assertion of \cref{thm:priced-deficit} therefore gives
\begin{equation}\label{eq:priced-scaled-deficit}
 d_{Z,m}(t)\longrightarrow qt^{3/4}
 \qquad\text{for every fixed }t>0,
\end{equation}
for every choice of the minimizing indices.

For $0<v<t$, the minimizers at the two endpoint prices give
\begin{equation}\label{eq:quantum-secants}
 -d_{Z,m}(t)\le
       \frac{G_{Z,m}(t)-G_{Z,m}(v)}{t-v}
       \le-d_{Z,m}(v).
\end{equation}
Indeed, evaluating the objective at $n_v$ gives the upper inequality,
and evaluating the objective at $n_t$ at price $v$ gives the lower one.
These inequalities do not require differentiability at a price with more
than one minimizing sector.

Take a fixed partition $v=t_0<t_1<\cdots<t_r=t$.  Summing
\cref{eq:quantum-secants} bounds $G_{Z,m}(t)-G_{Z,m}(v)$ by the two endpoint
Riemann sums for the negative deficits.  First use
\cref{eq:priced-scaled-deficit} at the finitely many partition points, and
then refine the partition.  Since $qt^{3/4}$ is continuous, this proves
\begin{equation}\label{eq:quantum-price-increments}
 G_{Z,m}(t)-G_{Z,m}(v)
       \longrightarrow-\frac47q\bigl(t^{7/4}-v^{7/4}\bigr).
\end{equation}
The behavior at zero price fixes the integration constant.  The neutral
sector is a competitor, whereas $E_{n_v}\ge E\ge E_Z(Z)-C$ by
\cref{eq:neutral-offset}; hence
\begin{equation}\label{eq:quantum-zero-endpoint}
 -Cm^{-7/3}-v\,d_{Z,m}(v)\le G_{Z,m}(v)\le0.
\end{equation}
For every fixed $v>0$ this gives
$\limsup|G_{Z,m}(v)|\le qv^{7/4}$.  Combining it with
\cref{eq:quantum-price-increments} and then sending $v\downarrow0$ yields
\begin{equation}\label{eq:quantum-price-energy}
 G_{Z,m}(t)\longrightarrow-\frac47qt^{7/4},\qquad
 \frac{E_{n_t}-E_Z(Z)}{m^{7/3}}
       =G_{Z,m}(t)+t\,d_{Z,m}(t)
       \longrightarrow\frac37qt^{7/4}.
\end{equation}
Thus no assertion about the convergence of the priced minimizers as
$t\downarrow0$ was needed.

Choose fixed prices
\[
 0<t_-<q^{-4/3}<t_+.
\]
By \cref{eq:priced-scaled-deficit}, their deficits are eventually smaller
and larger than $m$, respectively.  Thus
$n_{t_-}>Z-m>n_{t_+}$.  Monotonicity of the sector energies gives
\[
 E_{n_{t_-}}-E_Z(Z)
       \le E_Z(Z-m)-E_Z(Z)
       \le E_{n_{t_+}}-E_Z(Z).
\]
Apply \cref{eq:quantum-price-energy} and let both prices approach
$q^{-4/3}$.  The two bounds converge to
$\frac37q(q^{-4/3})^{7/4}=\frac37q^{-4/3}$, proving
\cref{eq:fixed-sector-limit}.
\end{proof}

\subsection{The same charge in Thomas--Fermi theory}

The classical weak-ionization law was established by B\'enilan and Brezis
\cite[Section~6, Corollary~5]{BenilanBrezis2004}, following the bounds and
chemical-potential integration argument of Lieb and Simon
\cite[Theorems~IV.11--IV.12]{LiebSimon1977}.
We give the argument in the present normalization to identify its
coefficient with the charge obtained from the conditional quantum limit.

We allow real nuclear charges and real electron masses in this subsection.
For clarity, the admissible finite-mass class can be written as
\[
 \mathcal X=\{\rho\ge0:\rho\in L^1(\R^3)\cap L^{5/3}(\R^3)\}.
\]
All terms of $\mathcal T_Z$ are finite on this class.  For the direct term,
interpolation gives $\rho\in L^{6/5}$, and the Sobolev inequality gives the
embedding $L^{6/5}\subset(\dot H^1)^*$ and the finite Coulomb energy.
For the nuclear term, use $|x|^{-1}\in L^{5/2}_{\mathrm{loc}}$ near zero
and the finite mass away from zero.  Thus this is precisely the finite-mass
class in the definition of $E_Z^{\TF}(M)$.

\begin{proposition}\label{prop:tf-identification}
With the two-spin coefficient $T$ fixed above,
\begin{equation}\label{eq:tf-unit-deficit}
 \lim_{Z\to\infty}
       \bigl[E_Z^{\TF}(Z-1)-E_Z^{\TF}(Z)\bigr]
       =\frac37q^{-4/3}.
\end{equation}
The limit holds through real charges $Z>1$.  Consequently
\cref{eq:tf-constant} holds for every fixed real $m>0$, with
\[
 a_{\TF}=\frac37q^{-4/3}>0.
\]
\end{proposition}

\begin{proof}
We identify the charge in the unit-price Thomas--Fermi problem, then
integrate in the price and use the exact density scaling.

\paragraph{Monotonicity in the electron mass.}
The energy $E_Z^{\TF}(M)$ is nonincreasing in $M$.  To see this without an
attainment assumption, truncate any admissible density of mass $M_1$ to a
large ball.  The kinetic, nuclear, and direct terms converge to the
original ones.  If $M_2\ge M_1$, place the missing mass in a nonnegative
smooth packet of diameter proportional to $R$, centered at distance
$R^2$ from that ball.  For a packet of mass $a$ and this spatial scale,
its kinetic and direct energies are $O(a^{5/3}R^{-2})$ and
$O(a^2R^{-1})$.  Its nuclear term and its interaction with the fixed
truncated density tend to zero as $R\to\infty$.  The mass $a$ stays
bounded during this construction.  Taking first $R\to\infty$, then the
truncation radius to infinity, and then the original energy down to its
infimum proves monotonicity.  The same argument shows that imposing mass
exactly $M$ or at most $M$ gives the same infimum.

\paragraph{The zero-price minimizer is neutral.}
Apply \cite[\FTFPair]{Foundation} to the single positive point
source $Z\delta_0$.  Their unconstrained problem on
$L^{5/3}\cap(\dot H^1)^*$ has a unique minimizer $\rho_0$, with finite
mass at most $2Z$.  Its field $\phi_0=V-K*\rho_0$ satisfies
\begin{equation}\label{eq:tf-neutral-field}
 \rho_0=k(\phi_0)_+^{3/2},\qquad
 \Delta\phi_0=-4\pi Z\delta_0+4\pi\rho_0,\qquad
 c\min\left(\frac Z{|x|},|x|^{-4}\right)
       \le\phi_0(x)\le C|x|^{-4}.
\end{equation}
The cited lemmas supply the lower bound, finite mass, and decay at infinity.
For completeness, the upper bound is independent of $Z$: on a ball
$B(x,R)$ with $R=|x|/2$, compare $\phi_0$ with
\[
 v(y)=\frac{C R^4}{(R^2-|y-x|^2)^4}.
\]
For sufficiently large universal $C$, direct differentiation gives
$\Delta v\le4\pi k v^{3/2}$.  The function diverges at the boundary,
and the positive-part comparison for
$\Delta\phi_0=4\pi k\phi_0^{3/2}$ in this ball gives
$\phi_0(x)\le CR^{-4}$, as asserted.

Uniqueness makes $\rho_0$ and $\phi_0$ radial.  Newton's spherical-mean
identity and finite mass give
\[
 Z-\int\rho_0=\lim_{d\to\infty}d\phi_0(d)=0,
\]
where the last equality follows from \cref{eq:tf-neutral-field}.
Therefore $\int\rho_0=Z$ and
\begin{equation}\label{eq:tf-neutral-minimum}
 E_Z^{\TF}(Z)=\min_{\rho\in\mathcal X}\mathcal T_Z(\rho).
\end{equation}
There is no discrepancy between the variational classes here: every
finite-mass admissible density belongs to the cited Coulomb-dual class,
and the cited minimizer has finite mass and finite energy terms, hence
belongs to $\mathcal X$.

\paragraph{The unit-price minimizer.}
Minimize $\mathcal T_Z(\rho)+\int\rho$ over $\mathcal X$.
Removing density from the region $V\le1$ cannot increase this functional:
the kinetic term, the contribution $\int(1-V)\rho$ there, and all removed
direct interactions are nonnegative.  We may therefore restrict to the
ball $B(0,Z)$.  Since $V\in L^{5/2}(B(0,Z))$, coercivity and weak lower
semicontinuity in $L^{5/3}$ give a minimizer, and strict convexity of the
kinetic term gives uniqueness.  Denote it by $\rho_Z$ and put
$\phi_Z=V-K*\rho_Z$.  Nonnegative variations, followed by two-sided
multiplicative variations on the positive density, give
\begin{equation}\label{eq:tf-priced-euler}
 \rho_Z=k\bigl((\phi_Z-1)_+\bigr)^{3/2},\qquad
 \Delta\phi_Z=-4\pi Z\delta_0+4\pi\rho_Z.
\end{equation}
This holds globally: outside $B(0,Z)$ one has $\phi_Z\le V\le1$.
The minimizer and its field are radial by uniqueness.

The following comparison is the unit-price case of the classical
chemical-potential comparison \cite[Lemma~5.1]{Solovej2003}; we include
the short maximum-principle proof. The two fields satisfy
\begin{equation}\label{eq:tf-price-comparison}
 \phi_0\le\phi_Z\le\phi_0+1.
\end{equation}
For the first inequality, on a positive set of
$\phi_0-\phi_Z-\eta$, with $\eta>0$, the difference of the two reaction
terms in \cref{eq:tf-neutral-field,eq:tf-priced-euler} is nonnegative.
Testing with that positive part gives a nonpositive squared-gradient
bound and hence a zero test.  The test is compactly supported because
both potentials tend to zero at infinity.  For the second inequality,
$\phi_0+1$ solves the same unit-price equation away from the common
point source, and the same comparison applies.  At the origin the point
singularities cancel; the remaining potentials are continuous and locally
$H^1$, with locally $L^{5/3}$ source densities.  Thus the compact tests
are justified by $H^1$ approximation.  Let $\eta\downarrow0$ to obtain
\cref{eq:tf-price-comparison}.  The continuity and decay used here hold for
Coulomb potentials of finite $L^1\cap L^{5/3}$ densities, as also proved
in \cite[\FTFExistence]{Foundation}.

Write
\[
 Q_Z=Z-\int\rho_Z.
\]
By \cref{eq:tf-price-comparison}, $\phi_Z\ge0$ and
$\rho_Z\le k\phi_0^{3/2}\le C|x|^{-6}$.  Newton's identity reads
\begin{equation}\label{eq:tf-priced-mean}
 \phi_Z(d)=\frac{Q_Z}{d}
       +\int_{|y|>d}\left(\frac1d-\frac1{|y|}\right)\rho_Z(y)\,dy.
\end{equation}
Finite mass gives $Q_Z=\lim_{d\to\infty}d\phi_Z(d)\ge0$.
At $d=1$, the nonnegative second term and
\cref{eq:tf-price-comparison} give $Q_Z\le\phi_Z(1)\le C+1$.
Using the exterior integral of $C|y|^{-6}$ in
\cref{eq:tf-priced-mean}, we conclude
\begin{equation}\label{eq:tf-uniform-exterior}
 0\le\phi_Z(d)\le\frac C d+Cd^{-4}.
\end{equation}
Choose a fixed sufficiently large $R$.  Then $\phi_Z<1$ on $d>R$ for
every $Z$, so \cref{eq:tf-priced-euler} implies $\rho_Z=0$ there and
$\phi_Z(d)=Q_Z/d$.

\paragraph{Identification of the limiting charge.}
Let $Z\to\infty$ through real values.  On every compact annulus,
\cref{eq:tf-neutral-field,eq:tf-price-comparison} bound $\phi_Z$ above
and below, and \cref{eq:tf-priced-euler} bounds its source.  Local Poisson
estimates give a locally uniformly convergent subsequence.  Its limit
$\phi_\infty$ satisfies \cref{eq:limit-pde} and
\[
 c|x|^{-4}\le\phi_\infty(x)\le1+C|x|^{-4}.
\]
The uniform estimate \cref{eq:tf-uniform-exterior} gives decay at infinity.
By the uniqueness in \cref{lem:singular-profile}, this is the same field
constructed in \cref{sec:limit}.  Its exterior charge is $q$.  At any
fixed radius beyond both exterior supports,
$Q_Z=d\phi_Z(d)\to d\phi_\infty(d)=q$.  Uniqueness of the subsequential
limit therefore proves
\begin{equation}\label{eq:tf-unit-charge}
 Q_Z\longrightarrow q\qquad(Z\longrightarrow\infty).
\end{equation}
The use of real charges is permitted by the arbitrary-positive-point-source
hypotheses of \cite[\FTFPair]{Foundation}.

\paragraph{Other prices and their energies.}
For $t>0$, let $\rho_{Z,t}$ minimize
$\mathcal T_Z(\rho)+t\int\rho$, and write
$Q_Z(t)=Z-\int\rho_{Z,t}$.  Existence and uniqueness follow as above,
now restricting to $B(0,Z/t)$.  With $b=t^{-1/4}$, the transformation
\[
 \widetilde\rho(x)=b^6\rho_{Z,t}(bx),\qquad
 \widetilde\phi(x)=b^4\bigl(V-K*\rho_{Z,t}\bigr)(bx)
\]
gives the unit-price minimizer and field at nuclear charge $b^3Z$.
Indeed, the transformed functional is exactly
\[
 \mathcal T_{b^3Z}(\widetilde\rho)+\int\widetilde\rho
       =b^7\left(\mathcal T_Z(\rho_{Z,t})
                                      +t\int\rho_{Z,t}\right),
\]
and this transformation is a bijection of admissible densities.
Its mass is multiplied by $b^3$, so
\begin{equation}\label{eq:tf-price-scaling}
 Q_Z(t)=t^{3/4}Q_{t^{-3/4}Z}
       \longrightarrow qt^{3/4}
       \qquad(t>0\text{ fixed}).
\end{equation}

Define
\[
 \Pi_Z(t)=\min_{\rho\in\mathcal X}
                  \left(\mathcal T_Z(\rho)+t\int\rho\right),\qquad
 G_Z^{\TF}(t)=\Pi_Z(t)-tZ-E_Z^{\TF}(Z).
\]
The endpoint-minimizer comparisons give, for $0<v<t$,
\[
 -Q_Z(t)\le
      \frac{G_Z^{\TF}(t)-G_Z^{\TF}(v)}{t-v}
      \le-Q_Z(v).
\]
The exact zero-price identity \cref{eq:tf-neutral-minimum} gives
\[
 -vQ_Z(v)\le G_Z^{\TF}(v)\le0.
\]
Consequently the same fixed-partition integration and then $v\downarrow0$
argument used in \cref{eq:quantum-price-increments,eq:quantum-zero-endpoint}
gives
\begin{equation}\label{eq:tf-priced-energy}
 G_Z^{\TF}(t)\longrightarrow-\frac47qt^{7/4},\qquad
 \mathcal T_Z(\rho_{Z,t})-E_Z^{\TF}(Z)
       \longrightarrow\frac37qt^{7/4}.
\end{equation}
Each priced minimizer is also a minimizer at its own mass: an improvement
at that mass would improve the priced objective.  For fixed
$t_-<q^{-4/3}<t_+$, \cref{eq:tf-price-scaling} places these two masses on
opposite sides of $Z-1$ for all large $Z$.  The mass monotonicity proved
above and \cref{eq:tf-priced-energy}, followed by
$t_-,t_+\to q^{-4/3}$, give \cref{eq:tf-unit-deficit}.

\paragraph{Every fixed removed mass.}
For $m>0$ the density scaling
\[
 \rho(x)=m^2\sigma(m^{1/3}x)
\]
multiplies the mass by $m$ and satisfies
$\mathcal T_Z(\rho)=m^{7/3}\mathcal T_{Z/m}(\sigma)$.
It is a bijection between the corresponding admissible classes.  Hence
\begin{equation}\label{eq:tf-mass-scaling}
 E_Z^{\TF}(N)=m^{7/3}E_{Z/m}^{\TF}(N/m),\qquad
 I_m^{\TF}(Z)=m^{7/3}I_1^{\TF}(Z/m).
\end{equation}
For fixed $m>0$, the real charge $Z/m$ tends to infinity with $Z$.
Applying \cref{eq:tf-unit-deficit} proves \cref{eq:tf-constant} and the
claimed value of $a_{\TF}$.  This also completes the coefficient assertion
of \cref{thm:priced-deficit}.
\end{proof}

\subsection{The specified order of limits}

\begin{proof}[Completion of the proof of \cref{thm:main}]
The joint conclusion \cref{eq:joint-limit} follows from
\cref{prop:fixed-sector-limit,prop:tf-identification}.
To deduce the iterated limits, suppose first that, along unbounded integers
$m_j$, the normalized inner upper limit exceeds $a_{\TF}$ by a fixed
$\eta>0$.  By the definition of a limsup, for each $j$ one may choose an
arbitrarily large integer $Z_j>m_j$, in particular $Z_j/m_j\ge j$, with
\[
 \frac{I_{m_j}(Z_j)}{m_j^{7/3}}>a_{\TF}+\eta/2.
\]
This contradicts \cref{eq:joint-limit}.  The same selection applies if
these inner upper limits are infinite.  Likewise, a normalized inner lower
limit less than $a_{\TF}-\eta$ along unbounded $m_j$ allows choosing
$Z_j/m_j\ge j$ with normalized energies below $a_{\TF}-\eta/2$, again a
contradiction.  Therefore
\[
 \limsup_{m\to\infty}
       \frac{\limsup_{Z\to\infty}I_m(Z)}{m^{7/3}}\le a_{\TF},
 \qquad
 \liminf_{m\to\infty}
       \frac{\liminf_{Z\to\infty}I_m(Z)}{m^{7/3}}\ge a_{\TF}.
\]
Since the inner lower limit never exceeds the inner upper limit, both
outer limits exist and equal $a_{\TF}$.  The construction always chooses
$Z$ only after $m$ is fixed; no convergence in $Z$ at a fixed $m$ has been
assumed.
\end{proof}
